\scanpage{101}
\mspage{16}
% archive sheet number: 100
% source lines 1--9
condition (i). Cependant, la démonstration que nous donnerons \inserted{au § 5} de (3.2) montre que \uline{si le th. de pureté est satisfait pour les} $(X,Y^{(i)})$ \uline{\unsure{qui ont} dimension} $\leq i$, \uline{alors la conclusion sur} $H^i(U_{\bar K},\unsure{\mathbb Z_\ell|Z_\ell})^P$ \uline{reste valable}. Or \deleted{\unsure{la}} cette condition est toujours satisfaite pour $i=1$, \unsure{où|car} elle résulte du th. de pureté \unsure{en dimension} $2$ de \inserted{Nagata--}Zariski (SGA 2 \ldots), \deleted{Nous} cf plus bas \ldots\ . Le raisonnement de 3.5 a donne donc \deleted{\unsure{également}} ceci :

% source lines 10--13
\deleted{\uline{\unsure{Corollaire}} \unsure{3.9}.} Soit $X_K$ une \deleted{\unsure{schéma lisse de}} \inserted{courbe algébrique} \deleted{lisse de dim. $\leq 1$ sur $K$.} lisse sur $K$, alors l'action du groupe d'inertie $\bar I$ sur $H^1(X_{\bar K},\mathbb Z_\ell)$ est \unsure{qsi.|qu.} unipotente \uline{d'échelon 2}.

% source line 14
Grâce à la réduction II 5.4, on en conclut

% source lines 15--19
\uline{Corollaire} 3.9. Soit $X_K$ un schéma \deleted{normal} de type fini sur $K$. Alors $E=H^1(X_{\bar K},\mathbb Z_\ell)$ est de t.f., \inserted{et} l'action du groupe d'inertie sur $E$ est \unsure{qsi|quasi} unipotente d'échelon 3, \deleted{\unsure{et}} et $=$ qsi unipotente d'échelon 2 si $X$ est normal.
